\subsection{What the CIF ranking rests on}
\label{sec:mechanism}

Of four targeted changes to the selected CIF configuration, only reducing the
forest to one tree lowers the score on every dataset, in both tasks
(\cref{tab:cif-ranking-ablation}). Split-count ranking, disabling bootstrap
sampling, and disabling feature muting change balanced accuracy by at most
0.001 on average, with intervals that include zero. In regression these three
changes lower the mean $R^2$, two of them with intervals that exclude zero, but
the three \nolinkurl{coepra} datasets set those means, and the medians stay
within 0.007 of zero.

\begin{table}[htb]
  \centering
  \footnotesize
  \setlength{\tabcolsep}{4pt}
  \renewcommand{\arraystretch}{1.08}
  \begin{tabular}{@{}llS[table-format=+1.3]S[table-format=+1.3]cc@{}}
    \toprule
    \tablehead{Task} & \tablehead{CIF change} & {Mean difference} & {Median difference} & 95\% interval & Wins--losses \\
    \midrule
    Classification & Split-count ranking & -0.001 & 0.000 & [$-$0.005, +0.002] & \winloss{12}{9} \\
    Classification & Disable bootstrap & -0.001 & 0.000 & [$-$0.004, +0.003] & \winloss{9}{12} \\
    Classification & Disable feature muting & +0.001 & 0.000 & [$-$0.001, +0.002] & \winloss{13}{5} \\
    Classification & Use one tree & -0.071 & -0.043 & [$-$0.101, $-$0.045] & \winloss{0}{22} \\
    \addlinespace
    Regression & Split-count ranking & -0.053 & -0.003 & [$-$0.142, $-$0.002] & \winloss{2}{6} \\
    Regression & Disable bootstrap & -0.064 & -0.007 & [$-$0.128, $-$0.010] & \winloss{4}{4} \\
    Regression & Disable feature muting & -0.157 & 0.000 & [$-$0.468, +0.001] & \winloss{4}{4} \\
    Regression & Use one tree & -0.658 & -0.059 & [$-$1.692, $-$0.054] & \winloss{0}{8} \\
    \bottomrule
  \end{tabular}
  \caption{CIF ablations on the real-data benchmark, 22 classification and 8
  regression datasets. Differences are the changed variant minus the selected
  CIF configuration in balanced accuracy or $R^2$, paired within seed and fold
  and averaged within each dataset over downstream learners, supported standard
  values of $k$, and the 25 paired replicates; positive values favor the
  variant. Intervals are percentile bootstrap intervals for the mean over
  datasets (20{,}000 replicates); wins and losses count datasets, with exact
  ties omitted.}
  \label{tab:cif-ranking-ablation}
\end{table}

\paragraph{Support.}
The support of a ranker is the set of features that win at least one split.
Features outside it receive zero importance and follow in the tie order of
\cref{sec:rank-construction}, so a ranking carries information only within its
support. On the synthetic datasets, random forests give every feature nonzero
importance on 15 of the 16 designs, and 800 of 1,000 on the 1,000-feature
classification design. On 13 of the 16 designs, CIF's support keeps 90 to
100\% of the planted features, so the support serves as a shortlist. A single
CIT splits on few features, so zero-importance filler makes up most of its
top-$k$ list: 52, 81, and 90\% of its top 10, 25, and 50 features, against 0,
3, and 14\% for CIF and none for random forests. This is consistent with CIT
ranking near the bottom of the benchmark.

\paragraph{Feature use in sampled forests.}
In sparse designs, feature sampling keeps informative features out of most CIF
splits (\cref{app:candidate-availability}). CIT, which considers every feature
at each node, uses informative features in 100.0\% of classification splits
and 93.5\% of regression splits. Over the 12 classification forest designs,
CIF uses them in 33.2\% of splits, against 91.7\% for CIF-all, 8.4\% for RF,
and 6.5\% for ExtraTrees. RF and ExtraTrees grow to purity, so most of their
splits are deep splits on noise, while CIF stops at the permutation test. With
32 of 1,000 features sampled at a node and two informative, an informative
feature is a candidate at 6.3\% of nodes, and when the sampled subset holds
none, Stage~A selects from the noise it was given.
